\documentclass[12pt,a4paper]{article}
\usepackage[utf8]{vietnam}
\usepackage{amsmath,amssymb}
\usepackage{geometry}
\geometry{top=1.5cm,bottom=1.2cm,left=1.4cm,right=1.4cm,headheight=28pt,headsep=12pt,footskip=18pt}
\usepackage{tikz}
\usepackage{xcolor}
\usepackage{enumerate}
\usepackage{graphicx}
\usepackage[version=4]{mhchem}
\usepackage{fancyhdr}
\usepackage{ex_test}

\renewcommand{\baselinestretch}{0.95}
\setlength{\parskip}{0pt}

\definecolor{hfRed}{RGB}{211,47,47}
\definecolor{hfGreen}{RGB}{139,195,144}

\pagestyle{fancy}
\fancyhf{}
\fancyhead[L]{\small\itshape\color{hfRed} Tài liệu lưu hành nội bộ}
\fancyhead[R]{\small\itshape\color{hfRed} Thầy \textbf{TRẦN VĂN THẠNH} - 0777.470.803}
\renewcommand{\headrulewidth}{0.8pt}
\renewcommand{\headrule}{\hbox to\headwidth{\color{hfGreen}\leaders\hrule height \headrulewidth\hfill}}

\fancyfoot[L]{\small\itshape\color{hfRed} CS1: 6/15 Nguyễn Hoàng, P. Kim Long, TP Huế \quad|\quad CS2: 24 Đặng Thái Thân, TP Huế}
\fancyfoot[R]{\small\bfseries\color{hfRed} Trang \thepage}
\renewcommand{\footrulewidth}{0.8pt}
\renewcommand{\footrule}{\hbox to\headwidth{\color{hfGreen}\leaders\hrule height \footrulewidth\hfill}}

\fancypagestyle{firstpage}{
  \fancyhf{}
  \renewcommand{\headrulewidth}{0pt}
  \fancyfoot[L]{\small\itshape\color{hfRed} CS1: 6/15 Nguyễn Hoàng, P. Kim Long, TP Huế \quad|\quad CS2: 24 Đặng Thái Thân, TP Huế}
  \fancyfoot[R]{\small\bfseries\color{hfRed} Trang \thepage}
  \renewcommand{\footrulewidth}{0.8pt}
  \renewcommand{\footrule}{\hbox to\headwidth{\color{hfGreen}\leaders\hrule height \footrulewidth\hfill}}
}

\begin{document}
\thispagestyle{firstpage}

\noindent
\begin{minipage}[t]{0.42\textwidth}
	\centering\fontsize{10pt}{12pt}\selectfont
	\textbf{THPT NGUYỄN ĐÌNH CHIỂU}\\
	\textbf{ĐỀ CHÍNH THỨC}\\
	\fbox{\textbf{MÃ ĐỀ: 1011}}
\end{minipage}\hfill
\begin{minipage}[t]{0.56\textwidth}
	\centering\small
	\textbf{ĐỀ KIỂM TRA GIỮA KÌ I NĂM HỌC 2025 - 2026}\\
	\textbf{Môn thi: HÓA HỌC 12}\\
	\textit{Thời gian làm bài: 45 phút (không tính thời gian phát đề)}\\
	\textbf{\color{hfRed}(BẢN GIÁO VIÊN -- CÓ LỜI GIẢI CHI TIẾT)}
\end{minipage}

\smallskip
\noindent\hrulefill
\smallskip

\noindent\textbf{Họ và tên học sinh:} \dotfill\ \textbf{Lớp: 12A1}

\medskip
\noindent
\textbf{A. PHẦN TRẮC NGHIỆM (7 điểm)}

\smallskip
\noindent
\textbf{PHẦN I (3 điểm -- 12 câu): Câu trắc nghiệm nhiều phương án lựa chọn.} \textit{Thí sinh trả lời từ câu 1 đến câu 12. Mỗi câu hỏi thí sinh chỉ chọn 1 phương án.}

\Opensolutionfile{ans}[ans-gv1011]

\begin{ex}
Cellulose trinitrate được điều chế từ cellulose và nitric acid đặc có xúc tác sulfuric acid đặc, nóng. Để có $44{,}55\text{ kg}$ cellulose trinitrate, cần dùng dung dịch chứa $m\text{ kg}$ nitric acid (hiệu suất phản ứng đạt 90\%). Giá trị của $m$ là
\choice
{$28{,}350\text{ kg}$}
{$21{,}234\text{ kg}$}
{$25{,}515\text{ kg}$}
{\True $31{,}500\text{ kg}$}
\loigiai{
Phương trình hóa học điều chế cellulose trinitrate:
\[\ce{[C6H7O2(OH)3]_n + 3n HNO3 ->[H2SO4 \text{ đặc}, t^\circ] [C6H7O2(ONO2)3]_n + 3n H2O}\]
Khối lượng mol một mắt xích:
$M_{\ce{[C6H7O2(ONO2)3]}} = 297\text{ g/mol}$; $3 M_{\ce{HNO3}} = 3 \times 63 = 189\text{ g/mol}$.\\
Theo lí thuyết, để tạo thành $44{,}55\text{ kg}$ cellulose trinitrate cần lượng $\ce{HNO3}$ là:
\[m_{\ce{HNO3 (LT)}} = \dfrac{44{,}55 \times 189}{297} = 28{,}35\text{ kg}.\]
Do hiệu suất của phản ứng đạt $90\%$, khối lượng $\ce{HNO3}$ thực tế cần dùng là:
\[m = \dfrac{28{,}35}{90\%} = 31{,}500\text{ kg}.\]
\textbf{Chọn D.}
}
\end{ex}

\begin{ex}
Methyl salicylate (chất X) là sản phẩm tự nhiên của rất nhiều loại cây, thường được kết hợp với các loại tinh dầu khác dùng làm thuốc bôi ngoài da, thuốc xoa bóp, cao dán giảm đau, chống viêm. X có công thức cấu tạo như sau:
\begin{center}
	\includegraphics[width=3.8cm]{San_Pham/images_crop/fig_cau2_methyl_salicylate.png}
\end{center}
Cho các phát biểu sau:
\begin{enumerate}[(a)]
	\item Công thức phân tử của X là $\ce{C8H8O3}$.
	\item Phân tử X chứa 31,58\% oxygen về khối lượng.
	\item $a\text{ mol X}$ phản ứng tối đa với $a\text{ mol Na}$, sinh ra $a\text{ mol H2}$.
	\item $a\text{ mol X}$ phản ứng tối đa với $2a\text{ mol NaOH}$.
	\item X là hợp chất hữu cơ tạp chức, chứa đồng thời chức ester và chức alcohol.
\end{enumerate}
Số phát biểu \textbf{sai} là
\choice
{1}
{\True 2}
{4}
{3}
\loigiai{
Phân tích từng phát biểu:
\begin{itemize}
	\item (a) Đúng. Công thức phân tử của X là $\ce{C8H8O3}$ (gồm 1 vòng benzen 6C, 1 nhóm $-\ce{COOCH3}$ và 1 nhóm $-\ce{OH}$).
	\item (b) Đúng. $\%O = \dfrac{3 \times 16}{152} \times 100\% \approx 31{,}58\%$.
	\item (c) Sai. Phân tử X chỉ có 1 nhóm $-\ce{OH}$ phenol có nguyên tử H linh động:
	\[\ce{X + Na -> X-Na + 1/2 H2 ^}\]
	Do đó $a\text{ mol X}$ tác dụng với $a\text{ mol Na}$ chỉ sinh ra $0{,}5a\text{ mol H2}$.
	\item (d) Đúng. X chứa 1 nhóm ester và 1 nhóm $-\ce{OH}$ phenol:
	\[\ce{HO-C6H4-COOCH3 + 2 NaOH -> NaO-C6H4-COONa + CH3OH + H2O}\]
	Vậy $a\text{ mol X}$ phản ứng tối đa với $2a\text{ mol NaOH}$.
	\item (e) Sai. X là hợp chất tạp chức chứa chức ester và chức \textbf{phenol} (nhóm $-\ce{OH}$ liên kết trực tiếp với vòng thơm), không phải chức alcohol.
\end{itemize}
Vậy có 2 phát biểu sai là (c) và (e).\\
\textbf{Chọn B.}
}
\end{ex}

\begin{ex}
Cho các chất: glucose, fructose, maltose, sucrose, cellulose. Số chất tham gia phản ứng với $\ce{Cu(OH)2}$ ở điều kiện thường tạo dung dịch có màu xanh lam là
\choice
{2}
{3}
{\True 4}
{5}
\loigiai{
Các chất có nhiều nhóm $-\ce{OH}$ kề nhau (polyalcohol) sẽ phản ứng với $\ce{Cu(OH)2}$ ở điều kiện thường tạo phức chất dung dịch màu xanh lam.\\
- Glucose, fructose, maltose, sucrose đều hòa tan $\ce{Cu(OH)2}$ ở điều kiện thường tạo dung dịch xanh lam (4 chất).\\
- Cellulose là polysaccharide có cấu trúc mạch dài kết tinh vững chắc, không hòa tan $\ce{Cu(OH)2}$ ở điều kiện thường.\\
\textbf{Chọn C.}
}
\end{ex}

\begin{ex}
Cho ester có công thức cấu tạo thu gọn $\ce{HCOOCH3}$. Tên gọi của ester này là
\choice
{\True methyl formate}
{formic acid}
{methyl acetate}
{ethyl formate}
\loigiai{
Ester $\ce{HCOOCH3}$ được tạo từ formic acid ($\ce{HCOOH}$) và methyl alcohol ($\ce{CH3OH}$). Tên gọi là methyl formate.\\
\textbf{Chọn A.}
}
\end{ex}

\begin{ex}
Cho các chất: sodium palmitate, sodium oleate, calcium stearate, potassium stearate. Số chất có thể là thành phần chính của xà phòng là
\choice
{\True 3}
{1}
{2}
{4}
\loigiai{
Xà phòng là hỗn hợp muối sodium hoặc potassium của các acid béo (thường là palmitic acid, stearic acid, oleic acid...).\\
- Các chất có thể là thành phần chính của xà phòng: sodium palmitate, sodium oleate, potassium stearate (3 chất).\\
- Calcium stearate là muối ít tan trong nước, không được dùng làm xà phòng.\\
\textbf{Chọn A.}
}
\end{ex}

\begin{ex}
Để rửa sạch ống nghiệm dính aniline, người ta có thể
\choice
{rửa bằng dung dịch $\ce{NaOH}$ rồi rửa lại bằng nước}
{\True rửa bằng dung dịch $\ce{HCl}$ rồi rửa lại bằng nước}
{rửa bằng nước xà phòng rồi rửa lại bằng nước}
{rửa bằng nước}
\loigiai{
Aniline ($\ce{C6H5NH2}$) là amine thơm ít tan trong nước, nhưng có tính base yếu nên phản ứng dễ dàng với dung dịch $\ce{HCl}$ tạo thành muối anilinium chloride ($\ce{C6H5NH3Cl}$) tan tốt trong nước:
\[\ce{C6H5NH2 + HCl -> C6H5NH3Cl}\]
Sau đó tráng lại bằng nước thì ống nghiệm sẽ sạch hoàn toàn.\\
\textbf{Chọn B.}
}
\end{ex}

\begin{ex}
Cho hợp chất hữu cơ X mạch hở có công thức cấu tạo như hình bên.
\centerline{\includegraphics[width=6.5cm]{San_Pham/images_crop/fig_cau7_malonate.png}}
Cho các phát biểu sau:
\begin{enumerate}[(1)]
	\item X là ester hai chức.
	\item Thủy phân X trong môi trường base tạo muối và alcohol.
	\item Công thức phân tử của X là $\ce{C7H12O4}$.
	\item Giữa các phân tử X có liên kết hydrogen liên phân tử.
\end{enumerate}
Số phát biểu \textbf{không đúng} là
\choice
{2}
{3}
{0}
{\True 1}
\loigiai{
Công thức cấu tạo của X là diethyl malonate: $\ce{CH2(COOC2H5)2}$ hay $\ce{C2H5OOC-CH2-COOC2H5}$.\\
- Phát biểu (1) đúng: X có hai nhóm chức ester $-\ce{COO}-$.\\
- Phát biểu (2) đúng: Thủy phân trong môi trường kiềm tạo muối sodium malonate $\ce{CH2(COONa)2}$ và ethanol $\ce{C2H5OH}$.\\
- Phát biểu (3) đúng: Công thức phân tử của X là $\ce{C7H12O4}$ ($M = 160\text{ g/mol}$).\\
- Phát biểu (4) sai: Phân tử ester không có nguyên tử H linh động nên không tạo được liên kết hydrogen liên phân tử giữa các phân tử X.\\
Vậy có 1 phát biểu không đúng là (4).\\
\textbf{Chọn D.}
}
\end{ex}

\begin{ex}
Khi thực hiện phản ứng xà phòng hóa chất béo bằng dung dịch $\ce{NaOH}$, người ta thường thêm dung dịch $\ce{NaCl}$ bão hòa vào hỗn hợp phản ứng chủ yếu để
\choice
{tăng tốc độ phản ứng thủy phân}
{giúp xà phòng tan tốt hơn trong nước}
{\True làm cho xà phòng nổi lên để dễ tách ra}
{tránh sự phân hủy của xà phòng ở nhiệt độ cao}
\loigiai{
Dung dịch $\ce{NaCl}$ bão hòa làm tăng khối lượng riêng của dung dịch và làm giảm độ tan của muối sodium của acid béo (hiện tượng kết muối -- salting out), giúp xà phòng nhẹ hơn nổi lên trên bề mặt, dễ dàng vớt và tách ra khỏi dung dịch.\\
\textbf{Chọn C.}
}
\end{ex}

\begin{ex}
Bảng sau đây ghi lại nhiệt độ nóng chảy của một số chất hữu cơ X, Y, Z, T ở điều kiện thường:
\begin{center}
\begin{tabular}{|c|c|c|c|c|}
\hline
\textbf{Chất} & \textbf{X} & \textbf{Y} & \textbf{Z} & \textbf{T} \\ \hline
\textbf{Nhiệt độ nóng chảy ($^\circ\text{C}$)} & $-114{,}1$ & $16{,}6$ & $-97{,}8$ & $205$ \\ \hline
\end{tabular}
\end{center}
Biết rằng T là glutamic acid. Nhận xét nào sau đây về T là \textbf{đúng}?
\choice
{Ở điều kiện thường, T là chất lỏng không màu, dễ bay hơi}
{\True Ở điều kiện thường, T tồn tại ở dạng ion lưỡng cực, có nhiệt độ nóng chảy cao và dễ tan trong nước}
{T chỉ có tính acid, không có tính base}
{T không có nhiệt độ nóng chảy xác định vì là chất vô định hình}
\loigiai{
Glutamic acid ($\ce{HOOC-CH2-CH2-CH(NH2)-COOH}$) là một amino acid. Ở điều kiện thường, amino acid tồn tại chủ yếu ở dạng ion lưỡng cực ($\ce{^-OOC-CH2-CH2-CH(NH3^+)-COO-}$), là chất rắn kết tinh, có nhiệt độ nóng chảy rất cao ($205^\circ\text{C}$) và tan tốt trong nước nhờ tương tác ion bền vững.\\
\textbf{Chọn B.}
}
\end{ex}

\begin{ex}
Cho công thức cấu tạo của chất hữu cơ X như hình bên:
\begin{center}
	\includegraphics[width=8.0cm]{San_Pham/images_crop/fig_cau10_maltose.png}
\end{center}
Tên gọi của chất X là
\choice
{glucose}
{\True maltose}
{sucrose}
{fructose}
\loigiai{
Hình ảnh biểu diễn 2 gốc $\alpha$-glucose liên kết qua liên kết $\alpha$-1,4-glycosidic, trong phân tử còn 1 nhóm $-\ce{OH}$ hemiacetal tự do. Đây là công thức cấu tạo của maltose (đường mạch nha).\\
\textbf{Chọn B.}
}
\end{ex}

\begin{ex}
Cho công thức cấu tạo của glutamic acid như sau:
\begin{center}
	\includegraphics[width=6.0cm]{San_Pham/images_crop/fig_cau11_glutamic_acid.png}
\end{center}
Cho các phát biểu sau:
\begin{enumerate}[(a)]
	\item Glutamic acid là một $\alpha$-amino acid.
	\item Muối monosodium glutamate được dùng làm mì chính (bột ngọt).
	\item Trong dung dịch $\text{pH} = 3{,}2$, glutamic acid tồn tại chủ yếu ở dạng anion tích điện âm.
	\item Glutamic acid có tính lưỡng tính.
\end{enumerate}
Số phát biểu đúng là
\choice
{1}
{2}
{4}
{\True 3}
\loigiai{
Phân tích từng phát biểu:
\begin{itemize}
	\item (a) Đúng. Nhóm amino $-\ce{NH2}$ đính vào carbon vị trí số 2 ($\alpha$).
	\item (b) Đúng. Monosodium glutamate ($\ce{HOOC-[CH2]2-CH(NH2)-COONa}$) chính là bột ngọt (mì chính).
	\item (c) Sai. Điểm đẳng điện của glutamic acid là $\text{pI} \approx 3{,}22$. Tại $\text{pH} = 3{,}2 \approx \text{pI}$, glutamic acid tồn tại chủ yếu ở dạng ion lưỡng cực trung hòa điện tích ròng, không phải dạng anion tích điện âm.
	\item (d) Đúng. Do chứa đồng thời nhóm $-\ce{COOH}$ và nhóm $-\ce{NH2}$ nên glutamic acid có tính lưỡng tính.
\end{itemize}
Vậy có 3 phát biểu đúng là (a), (b), (d).\\
\textbf{Chọn D.}
}
\end{ex}

\begin{ex}
Cho các phát biểu sau:
\begin{enumerate}[(a)]
	\item Dung dịch methylamine làm quỳ tím hóa xanh.
	\item Aniline là chất lỏng, không màu, ít tan trong nước, để lâu trong không khí bị oxy hóa chuyển sang màu đen nâu.
	\item Khí dimethylamine có mùi khai khó chịu, độc.
	\item Tính base của aniline mạnh hơn methylamine.
\end{enumerate}
Số phát biểu đúng là
\choice
{1}
{\True 3}
{2}
{4}
\loigiai{
Phân tích từng phát biểu:
\begin{itemize}
	\item (a) Đúng. Methylamine ($\ce{CH3NH2}$) có tính base mạnh hơn ammonia, dung dịch làm xanh quỳ tím.
	\item (b) Đúng. Aniline tinh khiết là chất lỏng không màu, để ngoài không khí dễ bị oxy hóa thành màu nâu đen.
	\item (c) Đúng. Dimethylamine ($\ce{(CH3)2NH}$) ở điều kiện thường là chất khí có mùi khai độc.
	\item (d) Sai. Do ảnh hưởng hút electron của gốc phenyl làm giảm mật độ e trên N, nên tính base của aniline yếu hơn methylamine (và yếu hơn $\ce{NH3}$).
\end{itemize}
Vậy có 3 phát biểu đúng là (a), (b), (c).\\
\textbf{Chọn B.}
}
\end{ex}

\Closesolutionfile{ans}

\vspace{0.3cm}
\noindent
\textbf{PHẦN II (3 điểm -- 3 câu): Câu trắc nghiệm đúng sai.} \textit{Thí sinh trả lời từ câu 1 đến câu 3. Trong mỗi ý a), b), c), d) ở mỗi câu thí sinh chọn đúng hoặc sai.}

\setcounter{ex}{0}

\begin{ex}
Triglyceride đóng vai trò là nguồn cung cấp năng lượng và chuyên chở các chất béo trong quá trình trao đổi chất. Cho triglyceride X có công thức cấu tạo như hình sau:
\begin{center}
	\includegraphics[width=13.5cm]{San_Pham/images_crop/fig_cau1_p2_triglyceride.png}
\end{center}
\begin{enumerate}[a)]
	\item X được cấu tạo từ các acid béo là palmitic acid, oleic acid và linoleic acid.
	\item Thủy phân hoàn toàn $427\text{ kg}$ triglyceride X bằng lượng $\ce{NaOH}$ dư thu được $441\text{ kg}$ muối của acid béo.
	\item Trong các acid béo cấu tạo nên X có một acid béo thuộc loại acid béo omega -8.
	\item Tổng số nguyên tử trong 1 phân tử X là 158.
\end{enumerate}
\loigiai{
\begin{itemize}
	\item \textbf{Ý a) Sai.} Quan sát công thức cấu tạo của 3 gốc acid béo trong X:
	\begin{itemize}
		\item Gốc 1: $\ce{C15H31COO-}$ (từ palmitic acid).
		\item Gốc 2: $\ce{CH3(CH2)7CH=CH(CH2)7COO-}$ tức $\ce{C17H33COO-}$ (từ oleic acid).
		\item Gốc 3: có 3 liên kết đôi $\ce{C=C}$, công thức là $\ce{CH3CH2(CH=CHCH2)3(CH2)6COO-}$ tức $\ce{C17H29COO-}$ (từ linolenic acid, chứ không phải linoleic acid có 2 liên kết đôi).
	\end{itemize}
	\item \textbf{Ý b) Đúng.} Công thức phân tử của X là $\ce{C55H98O6}$ ($M_X = 55 \times 12 + 98 + 6 \times 16 = 854\text{ g/mol}$).\\
	Số mol X: $n_X = \dfrac{427}{854} = 0{,}5\text{ kmol}$.\\
	Hỗn hợp muối thu được gồm:
	\[\ce{C15H31COONa}\ (278\text{ g/mol}) + \ce{C17H33COONa}\ (304\text{ g/mol}) + \ce{C17H29COONa}\ (300\text{ g/mol})\]
	Tổng khối lượng mol hỗn hợp muối $= 278 + 304 + 300 = 882\text{ g/mol}$.\\
	Khối lượng muối thu được: $m_{\text{muối}} = 0{,}5 \times 882 = 441\text{ kg}$.
	\item \textbf{Ý c) Sai.} Đánh số từ nguyên tử C ở nhóm methyl ($\ce{CH3-}$) xa nhất:\\
	- Gốc oleic acid có liên kết đôi đầu tiên ở vị trí C9 $\rightarrow$ là acid béo omega-9.\\
	- Gốc linolenic acid có liên kết đôi đầu tiên ở vị trí C3 $\rightarrow$ là acid béo omega-3.\\
	Không có acid béo nào thuộc loại omega-8.
	\item \textbf{Ý d) Sai.} Phân tử X có công thức $\ce{C55H98O6}$. Tổng số nguyên tử trong một phân tử X là:
	\[55 + 98 + 6 = 159 \neq 158.\]
\end{itemize}
\textbf{Đáp án:} a - S, b - Đ, c - S, d - S.
}
\end{ex}

\begin{ex}
Các peptide có phản ứng thủy phân trong môi trường acid và môi trường kiềm, ngoài ra các peptide có từ 2 liên kết peptide trở lên phản ứng với $\ce{Cu(OH)2}$ trong môi trường kiềm tạo thành phức chất màu tím đặc trưng, gọi là phản ứng màu biuret.
\begin{enumerate}[a)]
	\item Dung dịch của các polypeptide hoà tan $\ce{Cu(OH)2}$ cho dung dịch có màu xanh tím.
	\item Thủy phân hoàn toàn $0{,}1\text{ mol}$ Glu--Ala--Lys cần vừa đủ $300\text{ mL}$ dung dịch $\text{KOH 1M}$.
	\item Thủy phân hoàn toàn $2{,}17\text{ gam}$ tripeptide mạch hở X (được tạo nên từ hai $\alpha$-amino acid có công thức dạng $\ce{H2NC_xH_yCOOH}$) bằng dung dịch $\ce{NaOH}$ dư, thu được $3{,}19\text{ gam}$ muối. Mặt khác, thủy phân hoàn toàn $3{,}255\text{ gam}$ X bằng dung dịch $\ce{HCl}$ dư, thu được $5{,}4375\text{ gam}$ muối.
	\item Gly--Ala không có phản ứng màu biuret với $\ce{Cu(OH)2}$.
\end{enumerate}
\loigiai{
\begin{itemize}
	\item \textbf{Ý a) Đúng.} Các polypeptide có nhiều liên kết peptide nên hòa tan $\ce{Cu(OH)2}$ trong môi trường kiềm tạo dung dịch phức màu tím (hoặc xanh tím) đặc trưng (phản ứng màu biuret).
	\item \textbf{Ý b) Sai.} Glu chứa 2 nhóm $-\ce{COOH}$, Ala chứa 1 nhóm $-\ce{COOH}$, Lys chứa 1 nhóm $-\ce{COOH}$ và 2 nhóm $-\ce{NH2}$.\\
	Do đó tripeptide Glu--Ala--Lys có 2 nhóm $-\ce{COOH}$ tự do và 2 liên kết peptide.\\
	Khi thủy phân hoàn toàn trong dung dịch KOH:
	\[\text{Glu--Ala--Lys} + 4\text{KOH} \rightarrow \text{Muối} + 2\ce{H2O}\]
	Số mol KOH cần phản ứng: $n_{\text{KOH}} = 4 \times 0{,}1 = 0{,}4\text{ mol}$.\\
	Thể tích dung dịch KOH 1M cần là: $V = \dfrac{0{,}4}{1} = 0{,}4\text{ lít} = 400\text{ mL} \neq 300\text{ mL}$.
	\item \textbf{Ý c) Đúng.}
	- Phản ứng với NaOH:\\
	Tripeptide X được tạo từ các amino acid dạng $\ce{H2N-R-COOH}$ (1 nhóm amino, 1 nhóm carboxyl).\\
	$\ce{X + 3 NaOH -> Muối + H2O}$\\
	Bảo toàn khối lượng: $m_X + m_{\ce{NaOH}} = m_{\text{muối}} + m_{\ce{H2O}}$\\
	$\Rightarrow 2{,}17 + 40 \times 3 n_X = 3{,}19 + 18 n_X \Rightarrow 102 n_X = 1{,}02 \Rightarrow n_X = 0{,}01\text{ mol}$.\\
	Khối lượng mol của X: $M_X = \dfrac{2{,}17}{0{,}01} = 217\text{ g/mol}$.\\
	- Phản ứng với HCl:\\
	Khi lấy $3{,}255\text{ gam}$ X thì $n_{X} = \dfrac{3{,}255}{217} = 0{,}015\text{ mol}$.\\
	$\ce{X + 2 H2O + 3 HCl -> Muối}$\\
	Khối lượng muối thu được:\\
	$m_{\text{muối}} = m_X + m_{\ce{H2O}} + m_{\ce{HCl}} = 3{,}255 + (2 \times 0{,}015 \times 18) + (3 \times 0{,}015 \times 36{,}5) = 5{,}4375\text{ gam}$.\\
	Giá trị hoàn toàn khớp với đề bài.
	\item \textbf{Ý d) Đúng.} Gly--Ala là dipeptide (chỉ có 1 liên kết peptide), điều kiện để có phản ứng màu biuret là phải có từ 2 liên kết peptide trở lên, do đó Gly--Ala không có phản ứng màu biuret.
\end{itemize}
\textbf{Đáp án:} a - Đ, b - S, c - Đ, d - Đ.
}
\end{ex}

\begin{ex}
Khi nấu rượu hoặc làm bánh mì từ bột bánh mì có chứa men (một loại nấm) và đường. Khi bột được để ở nơi ấm áp, các tế bào nấm men sẽ ăn đường để lấy năng lượng. Enzyme trong nấm men xúc tác cho phản ứng được gọi là quá trình lên men sinh ra ethanol và khí carbon dioxide:
\[\ce{(C6H10O5)_n} \xrightarrow[\text{enzyme}]{+\ce{H2O}} \ce{n C6H12O6} \xrightarrow{\text{enzyme}} \ce{2n C2H5OH + 2n CO2}\]
\begin{enumerate}[a)]
	\item Để điều chế $5\text{ lít}$ ethyl alcohol $46^\circ$ cần $5{,}4\text{ kg}$ bột bánh mì trên (chứa 75\% tinh bột, còn lại là tạp chất trơ). Biết hiệu suất của cả quá trình là 80\% và khối lượng riêng của ethyl alcohol nguyên chất là $0{,}8\text{ g/mL}$.
	\item Nếu đun nóng hỗn hợp lên men để tách $\ce{C2H5OH}$, đó là phương pháp chưng cất dựa trên sự khác nhau về độ tan.
	\item Lên men rượu luôn xảy ra tối ưu ở khoảng $30 - 35^\circ\text{C}$, nếu nhiệt độ thấp nấm men sẽ chết.
	\item Nếu cung cấp nhiều oxy cho môi trường lên men, hiệu suất tạo ethanol sẽ tăng.
\end{enumerate}
\loigiai{
\begin{itemize}
	\item \textbf{Ý a) Đúng.}
	Thể tích ethanol nguyên chất: $V_{\ce{C2H5OH}} = 5 \times 46\% = 2{,}3\text{ lít} = 2300\text{ mL}$.\\
	Khối lượng ethanol: $m_{\ce{C2H5OH}} = 2300 \times 0{,}8 = 1840\text{ gam} \Rightarrow n_{\ce{C2H5OH}} = \dfrac{1840}{46} = 40\text{ mol}$.\\
	Sơ đồ phản ứng: $\ce{C6H10O5 -> 2 C2H5OH}$.\\
	Số mol tinh bột theo lý thuyết: $n_{\text{tinh bột (LT)}} = \dfrac{40}{2} = 20\text{ mol}$.\\
	Với hiệu suất toàn quá trình là $80\%$, số mol tinh bột thực tế cần:
	\[n_{\text{tinh bột (TT)}} = \dfrac{20}{80\%} = 25\text{ mol}.\]
	Khối lượng tinh bột nguyên chất: $m_{\text{tinh bột}} = 25 \times 162 = 4050\text{ gam} = 4{,}05\text{ kg}$.\\
	Khối lượng bột bánh mì (chứa 75\% tinh bột):
	\[m_{\text{bột bánh mì}} = \dfrac{4{,}05}{75\%} = 5{,}4\text{ kg}.\]
	\item \textbf{Ý b) Sai.} Phương pháp chưng cất dựa trên sự khác nhau về \textbf{nhiệt độ sôi} của các chất trong hỗn hợp lỏng (ethanol sôi ở $78{,}3^\circ\text{C}$, nước sôi ở $100^\circ\text{C}$), không phải dựa vào độ tan.
	\item \textbf{Ý c) Sai.} Ở nhiệt độ thấp, nấm men không bị chết mà chỉ bị giảm hoặc ức chế hoạt tính xúc tác của enzyme. Khi nhiệt độ tăng trở lại phù hợp, nấm men vẫn hoạt động bình thường.
	\item \textbf{Ý d) Sai.} Quá trình lên men rượu là quá trình yếm khí (kị khí). Nếu cung cấp nhiều không khí (oxygen), nấm men sẽ chuyển sang hô hấp hiếu khí, oxy hóa phân tử đường thành $\ce{CO2}$ và $\ce{H2O}$ làm giảm hiệu suất tạo ethanol.
\end{itemize}
\textbf{Đáp án:} a - Đ, b - S, c - S, d - S.
}
\end{ex}

\vspace{0.3cm}
\noindent
\textbf{PHẦN III (1 điểm -- 4 câu): Câu trắc nghiệm yêu cầu trả lời ngắn.} \textit{Thí sinh trả lời từ câu 1 đến câu 4.}

\setcounter{ex}{0}

\begin{ex}
Một loại chất béo có chứa 75\% tristearine về khối lượng. Xà phòng hóa hoàn toàn $17{,}8\text{ kg}$ chất béo này trong dung dịch $\text{KOH}$, đun nóng thu được $x$ bánh xà phòng. Biết rằng trong mỗi bánh xà phòng có chứa $63\text{ gam}$ potassium stearate. Giá trị của $x$ là bao nhiêu?
\loigiai{
Khối lượng tristearin trong $17{,}8\text{ kg}$ chất béo là:
\[m_{\text{tristearin}} = 17{,}8 \times 75\% = 13{,}35\text{ kg} = 13350\text{ gam}.\]
Số mol tristearin:
\[n_{\text{tristearin}} = \dfrac{13350}{890} = 15\text{ mol}.\]
Phương trình hóa học xà phòng hóa:
\[\ce{(C17H35COO)3C3H5 + 3 KOH -> 3 C17H35COOK + C3H5(OH)3}\]
Số mol potassium stearate thu được:
\[n_{\ce{C17H35COOK}} = 3 \times n_{\text{tristearin}} = 3 \times 15 = 45\text{ mol}.\]
Khối lượng potassium stearate thu được:
\[m_{\ce{C17H35COOK}} = 45 \times 322 = 14490\text{ gam}.\]
Số bánh xà phòng thu được:
\[x = \dfrac{14490}{63} = \mathbf{230\text{ bánh}}.\]
\textbf{Đáp số:} $\mathbf{230}$.
}
\end{ex}

\begin{ex}
Aspirin là một hợp chất được sử dụng làm giảm đau, hạ sốt được điều chế theo phản ứng sau:
\[\ce{(CH3CO)2O + HOC6H4COOH <=> CH3COOC6H4COOH + CH3COOH}\]
Để sản xuất 500 viên thuốc aspirin cần tối thiểu $20{,}7\text{ gam}$ salicylic acid. Biết rằng mỗi viên thuốc có chứa $y\text{ mg}$ aspirin và hiệu suất phản ứng đạt 75\%. Giá trị của $y$ là bao nhiêu?
\loigiai{
Số mol salicylic acid ($\ce{HOC6H4COOH}$, $M = 138\text{ g/mol}$) là:
\[n_{\text{acid}} = \dfrac{20{,}7}{138} = 0{,}15\text{ mol}.\]
Theo phương trình hóa học, tỉ lệ phản ứng là $1 : 1$. Với hiệu suất phản ứng $H = 75\%$, số mol aspirin ($\ce{CH3COOC6H4COOH}$, $M = 180\text{ g/mol}$) thu được là:
\[n_{\text{aspirin}} = 0{,}15 \times 75\% = 0{,}1125\text{ mol}.\]
Tổng khối lượng aspirin thu được:
\[m_{\text{aspirin}} = 0{,}1125 \times 180 = 20{,}25\text{ gam} = 20250\text{ mg}.\]
Lượng aspirin có trong mỗi viên thuốc:
\[y = \dfrac{20250}{500} = \mathbf{40{,}5\text{ mg}}.\]
\textbf{Đáp số:} $\mathbf{40{,}5}$.
}
\end{ex}

\begin{ex}
Cho các nhận định sau:
\begin{enumerate}[(a)]
	\item Protein dạng hình cầu không tan trong nước, còn dạng hình sợi tan trong nước.
	\item Một trong những tính chất hóa học đặc trưng của protein là phản ứng thủy phân.
	\item Phản ứng của protein với nitrous acid cho sản phẩm màu vàng.
	\item Sự đông tụ không làm thay đổi cấu tạo ban đầu của protein bị biến đổi.
	\item Trong cơ thể, enzyme không đóng vai trò là chất xúc tác sinh học.
\end{enumerate}
Số phát biểu đúng là
\loigiai{
Phân tích các nhận định:
\begin{itemize}
	\item (a) Sai. Protein dạng hình cầu tan trong nước tạo dung dịch keo (như albumin), trong khi protein hình sợi (như keratin của tóc, móng; collagen của da, gân) không tan trong nước.
	\item (b) Đúng. Phản ứng thủy phân (dưới tác dụng của acid, base hoặc enzyme) tạo thành các amino acid là phản ứng hóa học quan trọng đặc trưng của protein.
	\item (c) Sai. Phản ứng màu xanthoproteic của protein là phản ứng với dung dịch \textbf{nitric acid} ($\ce{HNO3}$) đậm đặc cho kết tủa vàng, không phải nitrous acid ($\ce{HNO2}$).
	\item (d) Sai. Sự đông tụ protein do nhiệt độ hay tác nhân hóa học làm biến tính và phá vỡ cấu trúc không gian tự nhiên ban đầu của protein.
	\item (e) Sai. Hầu hết các enzyme có bản chất là protein và đóng vai trò là chất xúc tác sinh học thiết yếu trong cơ thể.
\end{itemize}
Vậy chỉ có 1 phát biểu đúng là (b).\\
\textbf{Đáp số:} $\mathbf{1}$.
}
\end{ex}

\begin{ex}
Để tráng một số lượng gương soi có diện tích bề mặt $35\text{ dm}^2$ với độ dày $0{,}08\text{ }\mu\text{m}$ người ta đun nóng dung dịch chứa $30{,}6\text{ gam}$ glucose với một lượng dung dịch $\ce{AgNO3}$ trong amoniac. Biết khối lượng riêng của silver là $10{,}49\text{ kg/L}$, hiệu suất phản ứng tráng gương là 80\% (tính theo glucose). Có tối đa bao nhiêu chiếc gương soi được sản xuất ra? \textit{(Kết quả làm tròn đến hàng đơn vị)}.
\loigiai{
1. Số mol glucose tham gia phản ứng:
\[n_{\text{glucose}} = \dfrac{30{,}6}{180} = 0{,}17\text{ mol}.\]
2. Phản ứng tráng bạc: $\ce{C6H12O6 -> 2 Ag}$.\\
Số mol Ag sinh ra theo thực tế (với $H = 80\%$):
\[n_{\ce{Ag}} = 0{,}17 \times 2 \times 80\% = 0{,}272\text{ mol}.\]
3. Khối lượng bạc tạo thành:
\[m_{\ce{Ag}} = 0{,}272 \times 108 = 29{,}376\text{ gam}.\]
4. Khối lượng riêng của Ag: $D = 10{,}49\text{ kg/L} = 10{,}49\text{ g/cm}^3$.\\
Tổng thể tích Ag thu được:
\[V_{\text{tổng}} = \dfrac{m_{\ce{Ag}}}{D} = \dfrac{29{,}376}{10{,}49} \approx 2{,}80038\text{ cm}^3.\]
5. Thể tích lớp mạ bạc trên một chiếc gương:
\begin{itemize}
	\item Diện tích gương: $S = 35\text{ dm}^2 = 3500\text{ cm}^2$.
	\item Độ dày lớp bạc: $h = 0{,}08\text{ }\mu\text{m} = 0{,}08 \times 10^{-4}\text{ cm} = 8 \times 10^{-6}\text{ cm}$.
	\item Thể tích bạc trên 1 gương:
	\[V_1 = S \times h = 3500 \times (8 \times 10^{-6}) = 0{,}028\text{ cm}^3.\]
\end{itemize}
6. Số lượng chiếc gương tối đa sản xuất được:
\[N = \dfrac{V_{\text{tổng}}}{V_1} = \dfrac{2{,}80038}{0{,}028} \approx 100{,}013 \Rightarrow \mathbf{100\text{ chiếc}}.\]
\textbf{Đáp số:} $\mathbf{100}$.
}
\end{ex}

\vspace{0.4cm}
\noindent
\textbf{B. PHẦN TỰ LUẬN (3,0 ĐIỂM)}

\setcounter{ex}{0}

\begin{ex}
\textbf{(1,5 điểm)} Viết công thức cấu tạo, gọi tên thay thế và gốc chức của amin bậc 2, bậc 3 có công thức phân tử $\ce{C4H11N}$?
\loigiai{
Độ bất bão hòa của phân tử: $k = \dfrac{2 \times 4 + 2 + 1 - 11}{2} = 0 \Rightarrow \ce{C4H11N}$ là amine no, đơn chức, mạch hở.\\
Các đồng phân amine bậc 2 và bậc 3 cùng tên gọi:
\begin{enumerate}[\bfseries 1.]
	\item \textbf{Amine bậc 2} (dạng $\ce{R-NH-R'}$): có 3 đồng phân cấu tạo:
	\begin{itemize}
		\item $\ce{CH3-NH-CH2-CH2-CH3}$:
		\begin{itemize}
			\item Tên thay thế: \textbf{N-methylpropan-1-amine}
			\item Tên gốc -- chức: \textbf{methylpropylamine}
		\end{itemize}
		\item $\ce{CH3-NH-CH(CH3)-CH3}$:
		\begin{itemize}
			\item Tên thay thế: \textbf{N-methylpropan-2-amine}
			\item Tên gốc -- chức: \textbf{isopropylmethylamine}
		\end{itemize}
		\item $\ce{CH3-CH2-NH-CH2-CH3}$:
		\begin{itemize}
			\item Tên thay thế: \textbf{N-ethylethanamine}
			\item Tên gốc -- chức: \textbf{diethylamine}
		\end{itemize}
	\end{itemize}
	\item \textbf{Amine bậc 3} (dạng $\ce{R-N(R')-R''}$): có 1 đồng phân cấu tạo:
	\begin{itemize}
		\item $\ce{(CH3)2N-CH2-CH3}$ (hay $\ce{CH3-N(CH3)-CH2-CH3}$):
		\begin{itemize}
			\item Tên thay thế: \textbf{N,N-dimethylethanamine}
			\item Tên gốc -- chức: \textbf{ethyldimethylamine}
		\end{itemize}
	\end{itemize}
\end{enumerate}
}
\end{ex}

\begin{ex}
\textbf{(1,5 điểm)}
\begin{enumerate}[a.]
	\item So sánh điểm giống và khác nhau giữa xà phòng và chất giặt rửa tổng hợp?
	\item Trình bày cơ chế giặt rửa của chúng?
\end{enumerate}
\loigiai{
\textbf{a. So sánh xà phòng và chất giặt rửa tổng hợp:}
\begin{itemize}
	\item \textbf{Giống nhau:}
	\begin{itemize}
		\item Đều là các chất hoạt động bề mặt có khả năng làm sạch các vết bẩn dầu mỡ bám trên bề mặt sợi vải, da, đồ vật,...
		\item Về cấu tạo phân tử, đều gồm hai phần:
		\begin{itemize}
			\item Phần đầu ưa nước (phân cực, mang điện tích âm hoặc có nhóm phân cực mạnh).
			\item Phần đuôi kị nước (ưa dầu mỡ, là chuỗi hydrocarbon dài không phân cực).
		\end{itemize}
	\end{itemize}
	\item \textbf{Khác nhau:}
	\begin{itemize}
		\item \textit{Về thành phần hóa học:}
		\begin{itemize}
			\item Xà phòng: là hỗn hợp muối sodium hoặc potassium của các acid béo cao (như palmitate, stearate, oleate...).
			\item Chất giặt rửa tổng hợp: thường là muối sodium của các alkyl sulfate ($\ce{RO-SO3Na}$) hoặc alkylbenzenesulfonate ($\ce{R-C6H4-SO3Na}$).
		\end{itemize}
		\item \textit{Về nguồn gốc sản xuất:}
		\begin{itemize}
			\item Xà phòng: được sản xuất bằng cách thủy phân chất béo tự nhiên (mỡ động vật, dầu thực vật) trong dung dịch kiềm.
			\item Chất giặt rửa tổng hợp: được tổng hợp hóa học từ các sản phẩm của quá trình chế biến dầu mỏ.
		\end{itemize}
		\item \textit{Về khả năng dùng trong nước cứng:}
		\begin{itemize}
			\item Xà phòng: tạo kết tủa với các ion $\ce{Ca^{2+}}, \ce{Mg^{2+}}$ trong nước cứng (tạo $\ce{(RCOO)2Ca}, \ce{(RCOO)2Mg}$ kết tủa), làm giảm bọt, mất tác dụng giặt rửa và bám bẩn làm ố vàng vải sợi.
			\item Chất giặt rửa tổng hợp: các muối của ion $\ce{Ca^{2+}}, \ce{Mg^{2+}}$ vẫn tan tốt trong nước, do đó vẫn duy trì tác dụng giặt rửa tốt trong nước cứng.
		\end{itemize}
		\item \textit{Về tính thân thiện môi trường:}
		\begin{itemize}
			\item Xà phòng: dễ bị vi sinh vật phân hủy sinh học, thân thiện với môi trường sinh thái.
			\item Chất giặt rửa tổng hợp: đặc biệt loại có mạch hydrocarbon phân nhánh rất khó bị vi sinh vật phân hủy sinh học, có thể gây ô nhiễm nguồn nước và tích tụ lâu dài.
		\end{itemize}
	\end{itemize}
\end{itemize}

\textbf{b. Cơ chế giặt rửa:}
\begin{enumerate}[1.]
	\item \textbf{Định hướng phân tử:} Khi hòa tan vào nước, các phân tử chất giặt rửa định hướng: phần đuôi hydrocarbon kị nước (ưa dầu) cắm sâu vào vết bẩn dầu mỡ bám trên bề mặt sợi vải, còn phần đầu phân cực ưa nước hướng ra phía môi trường nước.
	\item \textbf{Nhũ tương hóa và bóc tách:} Dưới tác động cơ học (vò xát, chà, khuấy của máy giặt), các phân tử chất giặt rửa kéo các phân tử dầu mỡ ra khỏi bề mặt vải và phân chia vết dầu mỡ thành những giọt siêu nhỏ lơ lửng trong nước (tạo thành các hạt micelle).
	\item \textbf{Phân tán và cuốn trôi:} Do các đầu ưa nước trên bề mặt hạt micelle mang cùng điện tích (thường mang điện tích âm) đẩy nhau nên chúng không thể kết tụ lại với nhau thành vết bẩn lớn, các hạt dầu mỡ phân tán đều trong nước và dễ dàng bị dòng nước cuốn trôi đi khi xả sạch.
\end{enumerate}
}
\end{ex}

\begin{center}
\textbf{--- HẾT ---}
\end{center}

\end{document}
